\documentclass[12pt,a4paper]{article}

% ---------------------------------------------------------------------------
% Packages
% ---------------------------------------------------------------------------
\usepackage[utf8]{inputenc}
\usepackage[T1]{fontenc}
\usepackage[margin=2.5cm]{geometry}
\usepackage{xcolor}
\usepackage{listings}
\usepackage{enumitem}
\usepackage{fancyhdr}
\usepackage{titlesec}
\usepackage[hidelinks]{hyperref}

%mine
\usepackage{url} % for breakable \path{} instead of \texttt{}
\urlstyle{tt}
% allow line breaks after hyphens in \path / \url
% needs \makeatletter…\makeatother around it because of the @:
\makeatletter
\g@addto@macro\UrlBreaks{\do\-}
\makeatother
% chemical formula
\usepackage{chemfig}
% caption sizing
\usepackage{caption}
\captionsetup{labelfont=bf} % bold figure label

% to force straight quotes ''' in \lst
\usepackage{upquote}
\usepackage{textcomp}

% ---------------------------------------------------------------------------
% Colours and listing styles
% ---------------------------------------------------------------------------
\definecolor{codebg}{RGB}{245,245,245}
\definecolor{codeframe}{RGB}{200,200,200}
\definecolor{cmdkey}{RGB}{0,80,160}
\definecolor{comment}{RGB}{100,120,100}
\definecolor{strclr}{RGB}{160,60,60}

% changing dots line in toc, lof dflt=4.5
% Make @ symbol a letter
\makeatletter
\renewcommand*\@dotsep{-0.5}
% Make @ symbol an "other".
\makeatother

\lstdefinestyle{shell}{
  upquote=true,
  language=bash,
  backgroundcolor=\color{codebg},
  basicstyle=\ttfamily\footnotesize,
  keywordstyle=\color{cmdkey}\bfseries,
  commentstyle=\color{comment}\itshape,
  stringstyle=\color{strclr},
  frame=single,
  rulecolor=\color{codeframe},
  breaklines=true,
  breakatwhitespace=false,
  postbreak=\mbox{\textcolor{codeframe}{$\hookrightarrow$}\space},
  showstringspaces=false,
  columns=fullflexible,
  keepspaces=true,
  xleftmargin=6pt,
  framexleftmargin=6pt,
}

\lstdefinestyle{output}{
  upquote=true,
  backgroundcolor=\color{white},
  basicstyle=\ttfamily\scriptsize,
  frame=single,
  rulecolor=\color{codeframe},
  breaklines=true,
  showstringspaces=false,
  columns=fullflexible,
  keepspaces=true,
  xleftmargin=6pt,
}

\lstdefinestyle{python}{
  upquote=true,
  language=Python,
  backgroundcolor=\color{codebg},
  basicstyle=\ttfamily\footnotesize,
  keywordstyle=\color{cmdkey}\bfseries,
  commentstyle=\color{comment}\itshape,
  stringstyle=\color{strclr},
  frame=single,
  rulecolor=\color{codeframe},
  breaklines=true,
  showstringspaces=false,
  columns=fullflexible,
  keepspaces=true,
  xleftmargin=6pt,
}
  
% inline code shortcut
\lstMakeShortInline[basicstyle=\ttfamily,columns=fullflexible]|

% ---------------------------------------------------------------------------
% Headers / footers
% ---------------------------------------------------------------------------
\pagestyle{fancy}
\fancyhf{}
\lhead{\small Py-ChemShell 25.0.6 on HIGGS}
\rhead{\small\thepage}
\lfoot{\small Py-ChemShell with DL\_POLY + ORCA}
\renewcommand{\headrulewidth}{0.4pt}

% Section spacing
\titlespacing*{\section}{0pt}{1.4em}{0.6em}

% Handy boxed "note" environment
\newcommand{\note}[1]{%
  \par\vspace{4pt}\noindent
  \fbox{\parbox{0.97\linewidth}{\small\textbf{Note.} #1}}%
  \par\vspace{4pt}}

% ---------------------------------------------------------------------------
% Title
% ---------------------------------------------------------------------------
\title{\textbf{Py-ChemShell 25.0.6 installation on HIGGS\\[2pt]
with DL\_POLY 5 and ORCA 6.1.0}}

\author{
\large A step-by-step guide for the philomath curious about QM/MM\\[4pt]
\normalsize\texttt{Version 0.98}\\[4pt]
\date{14\textsuperscript{th} July 2026}
William \textsc{Houppy}\\[4pt]
\normalsize\url{https://github.com/biochemix/pychemshell-dlpoly-orca}}

\begin{document}
\maketitle
\thispagestyle{fancy}

\begin{abstract}
\noindent
This guide walks you through installing \textbf{Py-ChemShell 25.0.6} in your own
account on the HIGGS high-performance computing (HPC) cluster, from the source
\emph{tarball}, and using it to run a hybrid quantum-mechanics/molecular-mechanics
(QM/MM) calculation. The quantum region is handled by \textbf{ORCA} and the
molecular-mechanics region by \textbf{DL\_POLY~5}. Every command is shown
explicitly, together with the output you should expect, so that you can follow
along even if this is your first time using an HPC cluster. No administrator
(\emph{root}) rights are required: everything is built in your home directory and
loaded through the \emph{environment modules} system.
\end{abstract}

\tableofcontents
\vspace{1em}

% Start section numbering at 0 (Background is section 0)
\setcounter{section}{-1}
% Start figure numbering at 0
\setcounter{figure}{-1}

% ===========================================================================
\section{Background: what are we installing, and why?}
% ===========================================================================
\setcounter{subsection}{-1}

Many interesting chemical problems, for example how an enzyme catalyses a
reaction, involve a small, chemically active region embedded in a very large
environment. Treating \emph{everything} with quantum mechanics is far too
expensive. The \textbf{QM/MM} approach solves this by describing only the reactive
part with quantum mechanics (QM) and the surrounding bulk with cheaper
molecular-mechanics (MM) force fields.

\textbf{ChemShell} is the ``glue'' that couples a QM program to an MM program and
drives the calculation (energies, geometry optimisations, and so on). In this
guide we use three pieces of software:

\begin{itemize}
  \item \textbf{Py-ChemShell 25.0.6}: the Python version of ChemShell, the
        driver that orchestrates the QM/MM calculation.
  \item \textbf{ORCA 6.1.0}: the QM engine (already installed on HIGGS by the
        administrators; we simply call it).
  \item \textbf{DL\_POLY~5}: the MM engine, downloaded and compiled
        automatically by the ChemShell installer.
\end{itemize}

% ===========================================================================
\subsection{The HIGGS cluster in one paragraph}

HIGGS is managed with \textbf{EasyBuild} and the \textbf{environment modules}
system. Software is not installed system-wide; instead you \emph{load} the pieces
you need for a session with the \texttt{module load} command, and they disappear again
when you log out or run \texttt{module purge}. Jobs are submitted to compute nodes
through the \textbf{SLURM} scheduler. The compute nodes each have 192 physical CPU
cores (two sockets of 96 cores, no hyper-threading). A short \texttt{cpudebug}
partition (maximum one hour) is available for quick tests, and we use it
throughout this guide.

\note{Throughout, \texttt{\$HOME} means your home directory (e.g.\
\texttt{/hpcfs/home/<your-username>}) and \texttt{\$} at the start of a line
is the shell prompt (you do \emph{not} type it).}

% ===========================================================================
\section{Obtaining and unpacking the tarball}
% ===========================================================================
\setcounter{subsection}{-1}

Py-ChemShell is freely available and can be downloaded from its website (\url{https://www.chemshell.org}) as a tarball.
A \emph{tarball} is a compressed archive (extension \texttt{.tar.gz})
containing the ChemShell source code. Place it in a dedicated folder in your home
directory. Here we assume it is already in \texttt{\$HOME/chemshell/}.\\

\begin{lstlisting}[style=shell]
# 1. Move into a working folder for ChemShell
cd $HOME/chemshell

# 2. Check the tarball is there (about 35 MB)
ls -lh chemsh-py-25.0.6.tar.gz

# 3. Unpack it
tar xzf chemsh-py-25.0.6.tar.gz

# 4. Enter the source directory that was created
cd chemsh-py-25.0.6
ls
\end{lstlisting}

You should see the top-level contents of the distribution:

\begin{lstlisting}[style=output]
AUTHORS  chemsh  COPYING  COPYING.LESSER  docs  INSTALL  README.md  setup  tests
\end{lstlisting}

The \texttt{setup} script is the installer; the \texttt{tests} folder
contains ready-made examples that we will reuse later.

\note{Always unpack into a \emph{clean} directory. If you have to start over,
delete the \texttt{chemsh-py-25.0.6} folder and unpack the tarball again,
rather than building on top of a half-finished attempt.}

% ===========================================================================
\section{Loading and checking the environment modules}
% ===========================================================================
\setcounter{subsection}{-1}

Before building anything we choose a consistent set of compilers and libraries.
We use the \textbf{\texttt{foss/2024a}} toolchain (Free and Open-Source
Software): GCC 13.3.0, OpenMPI 5.0.3, and the OpenBLAS/FlexiBLAS, FFTW and
ScaLAPACK maths libraries. ChemShell also needs Python $\geq 3.10$, which is a
separate module.

% ===========================================================================
\subsection{Start from a clean environment}

\noindent % avoid page break inside
\begin{minipage}{\linewidth}
\begin{lstlisting}[style=shell]
module purge
module load foss/2024a
module list
\end{lstlisting}
\end{minipage}

\texttt{module list} should show the toolchain components, including
\texttt{GCC/13.3.0}, \path{OpenMPI/5.0.3}, \texttt{OpenBLAS/0.3.27},
\texttt{FlexiBLAS/3.4.4}, \texttt{FFTW/3.3.10} and
\texttt{ScaLAPACK/2.2.0-fb}.

% ===========================================================================
\subsection{Add a suitable Python}

The Python built for the \emph{same} compiler core (\texttt{GCCcore-13.3.0})
keeps everything consistent:\\

\noindent % avoid page break inside
\begin{minipage}{\linewidth}
\begin{lstlisting}[style=shell]
module load Python/3.12.3-GCCcore-13.3.0
which python3
python3 --version        # must report 3.12.3, NOT the system 3.9
\end{lstlisting}
\end{minipage}

\note{Loading Python may print an \emph{Lmod} warning that a few small helper
modules (\texttt{XZ}, \texttt{binutils}, \texttt{zlib}) were reloaded to
their \texttt{GCCcore-13.3.0} versions. This is harmless: the module system
is simply keeping the compiler chain consistent.}

% =========================================================================== 
\subsection{Locate ORCA}

ORCA is pre-installed. Find its module and the full path to the executable:

\begin{lstlisting}[style=shell]
module avail orca
# Expected: ORCA/6.1.0-gompi-2023b-avx2
\end{lstlisting}

The important detail for later: this ORCA was built with a \emph{different}
toolchain (\path{gompi-2023b}, i.e.\ OpenMPI 4.1.6), not the OpenMPI 5.0.3 we
build ChemShell against. We will handle that carefully in Section~\ref{sec:parallel}.
For now, simply record the executable's absolute path:

\begin{lstlisting}[style=shell]
module load ORCA/6.1.0-gompi-2023b-avx2
which orca
# /hpcfs/software/EasyBuild/easybuild/software/ORCA/6.1.0-gompi-2023b-avx2/bin/orca
module unload ORCA/6.1.0-gompi-2023b-avx2   # unload again before building ChemShell
\end{lstlisting}

% =========================================================================== 
\subsection{Check the compilers and library paths}

EasyBuild sets convenient variables (\texttt{\$EBROOT...}) pointing at each
library. Confirm the MPI compiler wrappers and the maths libraries resolve:

\begin{lstlisting}[style=shell]
which mpif90 mpicc mpicxx
echo $EBROOTFLEXIBLAS
echo $EBROOTFFTW
echo $EBROOTSCALAPACK
ls $EBROOTFLEXIBLAS/lib*/libflexiblas.so
\end{lstlisting}

The \texttt{mpif90}/\texttt{mpicc}/\texttt{mpicxx} wrappers should live
under \path{OpenMPI/5.0.3-GCC-13.3.0/bin}, and
\texttt{libflexiblas.so} should exist (it is present under both
\texttt{lib} and \texttt{lib64}; we use the \texttt{lib} copy).

% ===========================================================================
\section{Patching the installer (one small fix)}
\label{sec:patch}
% ===========================================================================
\setcounter{subsection}{-1}

The installer's default compilers are Intel (\texttt{ifort}/\texttt{icx}).
We are using GCC, and we must also enable MPI. On this version there is a small
bug in the installer: when MPI is switched on, one line looks up a dictionary with
a text key instead of a true/false value, which crashes the script with a
\texttt{KeyError}.

The fix is a one-character-level change: wrap the lookup in \texttt{bool(...)}.
A robust way to apply it (independent of the exact line number) is:

\noindent % avoid page break inside
\begin{minipage}{\linewidth}
\begin{lstlisting}[style=shell]
# Show the offending line first (optional)
grep -n "bools\[args.mpi\]" setup

# Apply the fix
sed -i "s/bools\[args.mpi\]/bools[bool(args.mpi)]/" setup

# Confirm it is now bools[bool(args.mpi)]
grep -n "bools\[bool(args.mpi)\]" setup
\end{lstlisting}
\end{minipage}

\note{This is a genuine upstream bug and should be reported to the ChemShell
developers. It does not affect the correctness of your results: it only
prevents the installer from running until patched.}

% ===========================================================================
\section{Configuring and building Py-ChemShell}
% ===========================================================================
\setcounter{subsection}{-1}

Now build ChemShell together with DL\_POLY~5. The installer downloads DL\_POLY
from its official repository and compiles it automatically. ORCA is \emph{not}
built here. It is an external program we call at run time, so it needs no
\path{--orca} flag.

Make sure \texttt{foss/2024a} and the Python module are loaded (and ORCA is
\emph{not} loaded), then run:

\begin{lstlisting}[style=shell]
./setup \
  -fc mpif90 -cc mpicc -cpc mpicxx \
  -mpi=openmpi \
  --python_root_dir $EBROOTPYTHON \
  --blas   $EBROOTFLEXIBLAS/lib/libflexiblas.so \
  --lapack $EBROOTFLEXIBLAS/lib/libflexiblas.so \
  --fftw   $EBROOTFFTW \
  --scalapack \
  --dl_poly \
  -j 4
\end{lstlisting}

What each part means:
\begin{itemize}[nosep]
  \item \texttt{-fc/-cc/-cpc} --- the Fortran, C and C++ MPI compiler wrappers.
  \item \texttt{-mpi=openmpi} --- build with MPI parallel support.
  \item \texttt{-\hspace{0mm}-python\_root\_dir \$EBROOTPYTHON} --- forces the correct Python
        3.12, not the old system Python.
  \item \texttt{-\hspace{0mm}-blas / -\hspace{0mm}-lapack} --- point at FlexiBLAS (our maths library).
  \item \texttt{-\hspace{0mm}-fftw / -\hspace{0mm}-scalapack} --- fast Fourier transforms and scalable
        linear algebra.
  \item \texttt{-\hspace{0mm}-dl\_poly} --- download and build the DL\_POLY~5 MM engine.
  \item \texttt{-j 4} --- compile using 4 parallel jobs (faster).
\end{itemize}

\phantom{i} % trick for blank line

The build takes a couple of minutes. Some warnings scroll past (for example
type-mismatch warnings in \texttt{dlf\_mpi.f90}, an \texttt{executable stack} note,
and two \texttt{not an ac\-ces\-si\-ble directory} warnings about the BLAS/LAPACK paths).
These are all \textbf{harmless}. The build has finished successfully when you see
progress reach 100\% and a line similar to:

\begin{lstlisting}[style=output]
[100%] Built target namd

 >>> Time used for building: 109.447 s
\end{lstlisting}

% ===========================================================================
\section{Verifying the build}
% ===========================================================================
\setcounter{subsection}{-1}

Check that the three key products exist and that the driver runs:

\begin{lstlisting}[style=shell]
ls -lh bin/gnu/chemsh.x bin/gnu/chemsh
ls -lh build/gnu/lib/libdl_poly.so
./bin/gnu/chemsh -v
\end{lstlisting}

The last command should print:

\begin{lstlisting}[style=output]
chemsh: ChemShell version 25.0
\end{lstlisting}

If you see the version string, Py-ChemShell is installed. \texttt{chemsh} is a
small launcher script; \texttt{chemsh.x} is the compiled program; and
\texttt{libdl\_poly.so} is the DL\_POLY interface library that ChemShell calls
in-process.

% ===========================================================================
\section{A first QM/MM test (serial)}
\label{sec:serial}
% ===========================================================================
\setcounter{subsection}{-1}

We start with a small, \emph{serial} (single-core) test to confirm the whole
chain (ChemShell $\rightarrow$ ORCA $\rightarrow$ DL\_POLY) works. Rather
than invent a system, we reuse a \emph{validated} example shipped in the
\texttt{tests} folder (an $n$-butanol molecule) and simply swap the QM engine
to ORCA. Because the example has a known reference energy, a correct result is
easy to recognise.

\begin{figure}[ht]
\centering
\setchemfig{atom sep=2.4em}
\chemfig{
  H-C(-[2]H)(-[6]H)-C(-[2]H)(-[6]H)-C(-[2]H)(-[6]H)%
  -[,,,,blue,line width=2.5pt]%
  \textcolor{red}{C}(-[2,,,,red]\textcolor{red}{H})(-[6,,,,red]\textcolor{red}{H})%
  -[,,,,red]\textcolor{red}{\charge{[extra sep=3pt]90=\:,270=\:}{O}}%
  -[,,,,red]\textcolor{red}{H}
}
\captionsetup{width=0.70\textwidth}
\caption{QM/MM partition of butanol-1.\\ The atoms and bonds are drawn in \textcolor{red}{\textbf{red}}
in the \textcolor{red}{\textbf{QM region}} (\texttt{qm\_region=range(5)}): the carbinol
carbon, its two hydrogens and the hydroxyl O--H, treated by ORCA. The propyl tail in \textbf{black} is the \textbf{MM region}, treated by DL\_POLY. The thick \textcolor{blue}{\textbf{blue}} bond joining the two carbons is the \textcolor{blue}{\textbf{QM/MM
boundary in mechanical embedding}} in which the interaction across the cut is handled
entirely by classical force field.}
\label{fig:butanol_qmmm}
\end{figure}

% ===========================================================================
\subsection{Set up a test directory}

\begin{lstlisting}[style=shell]
mkdir -p $HOME/chemshell/qmmm_test
cd $HOME/chemshell/qmmm_test

# Copy the shipped geometry and force field
cp $HOME/chemshell/chemsh-py-25.0.6/tests/qmmm/butanol.cjson .
cp $HOME/chemshell/chemsh-py-25.0.6/tests/qmmm/nwchem-dl_poly/butanol.ff .
\end{lstlisting}

% =========================================================================== 
\subsection{Write the ChemShell input}

Create the file \texttt{butanol\_orca\_dlpoly.py} by pasting this whole block into
the terminal:

\noindent % avoid page break inside
\begin{minipage}{\linewidth}
\begin{lstlisting}[style=shell]
cat > butanol_orca_dlpoly.py << 'EOF'
from chemsh import *

# Load the molecule (geometry supplied with ChemShell)
frag = Fragment(coords='butanol.cjson')

# QM engine: ORCA, Hartree-Fock with a small basis, run on 1 core
qm = ORCA(method="hf", basis="6-31g", nprocs=1)

# MM engine: DL_POLY, using the shipped force field
mm = DL_POLY(ff='butanol.ff', rcut=100.0)

# Couple them: first 5 atoms are the QM region, mechanical embedding
qmmm = QMMM(frag=frag, qm_region=range(5), qm=qm, mm=mm, embedding='mechanical')

# One geometry optimisation
opt = Opt(theory=qmmm, algorithm="lbfgs", tolerance=0.0003, 
          maxcycle=300, maxene=400, trust_radius='energy')
opt.run()

print("QM/MM optimised energy =", opt.result.energy, "Hartree")
EOF
\end{lstlisting}
\end{minipage}

% =========================================================================== 
\subsection{Write the SLURM job script}

Because ORCA is run on a single core here, there is no MPI to worry about yet: we only need ORCA on the \texttt{PATH} and its own library directory on \texttt{LD\_LIBRARY\_PATH}. Create \texttt{run\_qmmm.slurm} by pasting this block into the terminal:\\

\noindent % avoid page break inside
\begin{minipage}{\linewidth}
\begin{lstlisting}[style=shell]
cat > run_qmmm.slurm << 'EOF'
#!/bin/bash --login
#SBATCH --job-name=qmmm_smoke
#SBATCH --partition=cpudebug
#SBATCH --nodes=1
#SBATCH --ntasks=1
#SBATCH --time=00:20:00
#SBATCH --output=qmmm_smoke.%j.out

# ChemShell build environment
module purge
module load foss/2024a Python/3.12.3-GCCcore-13.3.0

# Point at ORCA (executable + its own libraries). Do NOT module-load ORCA here.
ORCADIR=/hpcfs/software/EasyBuild/easybuild/software/ORCA/6.1.0-gompi-2023b-avx2
export PATH=$ORCADIR/bin:$PATH
export LD_LIBRARY_PATH=$ORCADIR/lib:$LD_LIBRARY_PATH

CHEMSH=$HOME/chemshell/chemsh-py-25.0.6/bin/gnu/chemsh

cd $SLURM_SUBMIT_DIR
$CHEMSH butanol_orca_dlpoly.py > butanol_orca_dlpoly.log 2>&1
echo "chemsh exit=$?"
EOF
\end{lstlisting}
\end{minipage}

% =========================================================================== 
\subsection{Submit and check}

\begin{lstlisting}[style=shell]
sbatch run_qmmm.slurm
squeue -u $USER            # watch it run (a minute or so)
\end{lstlisting}

When it finishes, inspect the results:

\begin{lstlisting}[style=shell]
cat qmmm_smoke.*.out
tail -5 butanol_orca_dlpoly.log
\end{lstlisting}

A successful run prints \texttt{chemsh exit=0} and an energy very close to the
ORCA reference value of $-114.99042281$~Hartree, for example:

\begin{lstlisting}[style=output]
=== QM/MM optimised energy = [-114.99042314] Hartree ===
 ChemShell exiting normally.
\end{lstlisting}

Agreement to within about $3\times10^{-7}$ Hartree confirms the full chain works.

% ===========================================================================
\section{Setting up Py-ChemShell with parallel ORCA (bu\-ta\-nol-1 example)}
\label{sec:parallel}
% ===========================================================================
\setcounter{subsection}{-1}

In real QM/MM work, the QM calculation dominates the run time (often around
95\%), so we want ORCA to use many CPU cores. This section works through a
complete parallel example on the same butanol-1 system used in
Section~\ref{sec:serial}, so you can compare the two directly. It introduces one
subtlety that must be handled correctly, then gives every file you need.

% =========================================================================== 
\subsection{Why a wrapper script is needed}

ChemShell was built against \textbf{OpenMPI 5.0.3}, but ORCA was built against
\textbf{OpenMPI 4.1.6}. Both provide a library with the \emph{same} file name
(\texttt{libmpi.so.40}). If both versions were visible at once, ChemShell could
accidentally pick up ORCA's 4.1.6 and misbehave in ways that are hard to
diagnose.

The clean solution is to give each program its own private environment. We keep
ChemShell in the \texttt{foss/2024a} environment, and we launch ORCA through a
small \textbf{wrapper script} that loads ORCA's own modules \emph{inside a
separate subprocess}. The two never see each other's MPI.

Create \texttt{orca\_parallel.sh} in the test directory and make it executable
by pasting this block into the terminal:

\begin{lstlisting}[style=shell]
cat > orca_parallel.sh << 'EOF'
#!/bin/bash --login
# Isolated launcher: ORCA gets its OWN OpenMPI 4.1.6 environment,
# separate from the ChemShell process that called it.
module purge
module load ORCA/6.1.0-gompi-2023b-avx2

# ORCA needs its full absolute path to spawn MPI workers correctly:
exec /hpcfs/software/EasyBuild/easybuild/software/ORCA/6.1.0-gompi-2023b-avx2/bin/orca "$@"
EOF
\end{lstlisting}

\begin{lstlisting}[style=shell]
chmod +x orca_parallel.sh
\end{lstlisting}

% =========================================================================== 
\subsection{Tell ChemShell to use the wrapper}

Create a parallel version of the input, \texttt{butanol\_orca\_dlpoly\_par.py}.
The only changes from the serial version are (i) \texttt{path} points at the
wrapper, and (ii) the number of ORCA ranks is read from SLURM automatically (see
Section~\ref{sec:onenumber}). Paste this block into the terminal:\\

\noindent % avoid page break inside
\begin{minipage}{\linewidth}
\begin{lstlisting}[style=shell]
cat > butanol_orca_dlpoly_par.py << 'EOF'
import os
from chemsh import *

# Number of ORCA MPI ranks, taken from the SLURM job (falls back to 1)
norca = int(os.environ.get('SLURM_NTASKS', 1))

frag = Fragment(coords='butanol.cjson')

qm = ORCA(method="hf", basis="6-31g",
          path="./orca_parallel.sh",   # launch ORCA through the wrapper
          nprocs=norca)                # run ORCA in parallel

mm = DL_POLY(ff='butanol.ff', rcut=100.0)

qmmm = QMMM(frag=frag, qm_region=range(5), qm=qm, mm=mm,
            embedding='mechanical')

opt = Opt(theory=qmmm, algorithm="lbfgs", tolerance=0.0003,
          maxcycle=300, maxene=400, trust_radius='energy')
opt.run()

print("QM/MM optimised energy =", opt.result.energy, "Hartree")
EOF
\end{lstlisting}
\end{minipage}

\note{In this setup ChemShell runs as a single process and calls DL\_POLY
in-process (through \texttt{libdl\_poly.so}), while only ORCA runs in parallel.
The QM and MM steps happen one after another, not at the same time, so you do
\emph{not} need to divide cores between ORCA and DL\_POLY: you simply size the
job for ORCA.}

% =========================================================================== 
\subsection{One setting for the core count: passing it via SLURM}

\label{sec:onenumber}

It is easy to make mistakes if the number of cores is typed in two places (the
job script and the ChemShell input). We avoid this by setting it \textbf{once} in
the SLURM script. SLURM automatically exports \path{--ntasks} as the
environment variable \texttt{\$SLURM\_NTASKS}, and our Python input reads that
variable (the \path{os.environ.get(}\texttt{\textquotesingle SLURM\_NTASKS\textquotesingle}\texttt{, 1)} line above).\\

Create the parallel job script \texttt{run\_qmmm\_par.slurm} by pasting this block
into the terminal:

\noindent % avoid page break inside
\begin{minipage}{\linewidth}
\begin{lstlisting}[style=shell]
cat > run_qmmm_par.slurm << 'EOF'
#!/bin/bash --login
#SBATCH --job-name=qmmm_par
#SBATCH --partition=cpudebug
#SBATCH --nodes=1
#SBATCH --ntasks=4               # <-- the ONLY place the core count is set
#SBATCH --time=00:20:00
#SBATCH --output=qmmm_par.%j.out

# ChemShell environment only (OpenMPI 5.0.3). ORCA's 4.1.6 is loaded
# separately, inside orca_parallel.sh.
module purge
module load foss/2024a Python/3.12.3-GCCcore-13.3.0

CHEMSH=$HOME/chemshell/chemsh-py-25.0.6/bin/gnu/chemsh

echo "ORCA ranks (SLURM_NTASKS) = $SLURM_NTASKS"
cd $SLURM_SUBMIT_DIR
$CHEMSH butanol_orca_dlpoly_par.py > butanol_orca_dlpoly_par.log 2>&1
echo "chemsh exit=$?"
EOF
\end{lstlisting}
\end{minipage}

To change the number of cores later, edit only the \path{--ntasks} line and
the value flows automatically into ORCA. Because \texttt{-\hspace{0mm}-nodes=1} places all
the tasks on a single node, \texttt{-\hspace{0mm}-ntasks=4} already reserves the four
``slots'' on that node for ORCA's internal \texttt{mpirun}, so no
\texttt{-\hspace{0mm}-ntasks-per-node} line is needed.

% =========================================================================== 
\subsection{Submit and check}

\begin{lstlisting}[style=shell]
sbatch run_qmmm_par.slurm
squeue -u $USER
\end{lstlisting}

After it finishes:

\begin{lstlisting}[style=shell]
cat qmmm_par.*.out
tail -5 butanol_orca_dlpoly_par.log
\end{lstlisting}

You should again get \texttt{chemsh exit=0} and the same energy
($\approx -114.99042$~Hartree) as the serial run: the correct answer, now with ORCA
running on several cores.

\note{Parallel run is far slower than serial on a tiny test. This is expected for
very small systems: parallel overhead dominates. Benchmark on a realistic-sized
molecule instead.}

\noindent % avoid page break inside
\begin{minipage}{\linewidth}

% =========================================================================== 
\subsection{Recap: the five files for the parallel butanol-1 example}

For reference, a complete parallel QM/MM run of butanol-1 uses these files, all
in \path{$HOME/chemshell/qmmm_test}:\\

\begin{itemize}
  \item \texttt{butanol.cjson}: the geometry, copied from the shipped tests
        (Section~\ref{sec:serial}).
  \item \texttt{butanol.ff}: the force field, copied from the shipped tests
        (Section~\ref{sec:serial}).
  \item \texttt{orca\_parallel.sh}: the wrapper that gives ORCA its own MPI
        environment.
  \item \texttt{butanol\_orca\_dlpoly\_par.py}: the ChemShell input, which calls
        ORCA through the wrapper and reads the rank count from
        \texttt{\$SLURM\_NTASKS}.
  \item \texttt{run\_qmmm\_par.slurm}: the job script, where \path{--ntasks}
        is the single place you choose how many cores ORCA uses.
\end{itemize}

To run more cores, change only \path{--ntasks} then restart with \texttt{sbatch run\_qmmm\_par.slurm}.

\end{minipage}

% ===========================================================================
\section{Summary}
% ===========================================================================
\setcounter{subsection}{-1}

You have:
\begin{enumerate}
  \item unpacked the Py-ChemShell 25.0.6 tarball into your home directory;
  \item loaded a consistent open-source toolchain (\texttt{foss/2024a}) plus a
        suitable Python;
  \item applied the one required fix to the installer;
  \item built Py-ChemShell together with DL\_POLY~5;
  \item verified the build and ran a validated serial QM/MM test with ORCA;
  \item set up parallel ORCA cleanly using an isolating wrapper script; and
  \item driven the core count from a single SLURM setting.
\end{enumerate}

Your installation uses only open-source components for ChemShell and DL\_POLY, and
calls the pre-installed ORCA as an external program: a robust, reproducible
setup for QM/MM chemistry on HIGGS.

\phantom{i} % trick for line space

It does not yet exercise electrostatic embedding, link-atom QM/MM boundaries, or periodic MM electrostatics (Ewald/PME), which a protein/water run requires. These are features of the installed stack rather than new software, but remain untested in this build. A small electrostatic-embedding and link-atom test would confirm those code paths before production enzyme calculations. Actual larger QM/MM experiments would have to be carefully configured and tuned with the optimal number of CPU cores.

\newpage

% ===========================================================================
\appendix
\section{Troubleshooting}
% ===========================================================================

\phantom{i} % trick for line space

\begin{itemize}[leftmargin=1.2em]
  \item \textbf{The installer stops with} \texttt{KeyError:\textquotesingle openmpi\textquotesingle} \textbf{during}         \texttt{./setup}\textbf{.}\\
        The installer was not patched. Apply the \texttt{sed} fix in
        Section~\ref{sec:patch} and run \texttt{./setup} again.

  \item \textbf{CMake reports} \texttt{Could NOT find Python3}
        \textbf{or a} \texttt{Wrong version} \textbf{message.}\\
        The Python module was not loaded, so the installer found the old system
        Python 3.9. Load \texttt{Python/3.12.3-GCCcore-13.3.0} and pass
        \texttt{-\hspace{0mm}-python\_root\_dir} \texttt{\$EBROOTPYTHON} to \texttt{./setup}.

  \item \textbf{The parallel ORCA job aborts during} \texttt{Startup}
        \textbf{with a} \texttt{not enough slots} \textbf{message.}\\
        SLURM did not reserve enough slots. With \path{--nodes=1}, make sure
        \path{--ntasks} equals the number of ORCA ranks you asked for (the
        \texttt{nprocs} value ORCA uses).

  \item \textbf{You want to start the build again from scratch.}\\
        Run \texttt{./setup -\hspace{0mm}-clean} in the source directory, or delete the
        \texttt{chemsh-py-25.0.6} folder and re-unpack the tarball.
\end{itemize}

\end{document}
